\documentclass[12pt]{article}

\usepackage[utf8]{inputenc}
\usepackage[T1]{fontenc}
\usepackage{lmodern}
\usepackage[spanish]{babel}
\usepackage{booktabs}
\usepackage{amsmath}
\usepackage{forest}
\usepackage{float}
\usepackage{listings}
\usepackage{xcolor}

\definecolor{codegreen}{rgb}{0,0.6,0}
\definecolor{codegray}{rgb}{0.5,0.5,0.5}
\definecolor{codepurple}{rgb}{0.58,0,0.82}
\definecolor{backcolour}{rgb}{0.95,0.95,0.92}

\lstdefinestyle{mystyle}{
    backgroundcolor=\color{backcolour},   
    commentstyle=\color{codegreen},
    keywordstyle=\color{magenta},
    numberstyle=\tiny\color{codegray},
    stringstyle=\color{codepurple},
    basicstyle=\ttfamily\footnotesize,
    breakatwhitespace=false,         
    breaklines=true,                 
    captionpos=b,                    
    keepspaces=true,                 
    numbers=left,                    
    numbersep=5pt,                  
    showspaces=false,                
    showstringspaces=false,
    showtabs=false,                  
    tabsize=2
}

\lstset{style=mystyle}

\sloppy
\setlength{\parindent}{0pt}

\usepackage[utf8]{inputenc}  
\usepackage[T1]{fontenc}     
\usepackage{lmodern}         

\begin{document}

% Título y materia
\begin{center}
  {\LARGE \textbf{Naive Bayes}}\\[0.5em]
  {\normalsize (Clasificador Bayesiano Ingenuo)}\\[0.5em]
  {Analítica de Datos, Universidad de San Andrés}
\end{center}

Si encuentran algún error en el documento o hay alguna duda, mandenmé un mail a rodriguezf@udesa.edu.ar y lo revisamos.

% Introducción a Naive Bayes
\section{Introducción a Naive Bayes}
El clasificador Naive Bayes es un algoritmo de aprendizaje automático basado en el teorema de Bayes, que asume la independencia condicional entre las características dado el valor de la variable objetivo. A pesar de ser una suposición ingenua (naive), el algoritmo puede ser muy efectivo en algunas situaciones.

\vspace{1em}

El teorema de Bayes establece que:
\[
P(Y|X) = \frac{P(X|Y)P(Y)}{P(X)}
\]
donde:
\begin{itemize}
    \item $P(Y|X)$ es la probabilidad posterior
    \item $P(X|Y)$ es la verosimilitud
    \item $P(Y)$ es la probabilidad a priori
    \item $P(X)$ es la evidencia
\end{itemize}

\vspace{1em}

En el contexto de clasificación, queremos encontrar la clase $Y$ que maximiza $P(Y|X)$ dado un conjunto de características $X$. Debido a la suposición de independencia, podemos escribir:
\[
P(X|Y) = P(X_1|Y) \times P(X_2|Y) \times \cdots \times P(X_n|Y)
\]

\section{Tipos de Naive Bayes}
Antes de aplicar un modelo Naive Bayes, es importante conocer las variantes que existen y cuándo conviene usar cada una, según la documentación oficial de \texttt{scikit-learn}:

\begin{itemize}
    \item \textbf{BernoulliNB:} Utiliza variables binarias (presencia o ausencia de características). Es ideal cuando los datos son de tipo binario, por ejemplo, presencia/ausencia de palabras en texto.
    \item \textbf{CategoricalNB:} Se utiliza para datos categóricos, donde las características son discretas y tienen más de dos posibles valores, común en datos no numéricos.
    \item \textbf{MultinomialNB:} Ideal para datos discretos como conteos de palabras en texto, especialmente cuando las características tienen más de dos valores posibles y están basadas en frecuencias.
    \item \textbf{ComplementNB:} Variante de MultinomialNB que mejora el rendimiento en conjuntos de datos desbalanceados, ajustando las probabilidades de las clases complementarias para compensar el sesgo.
    \item \textbf{GaussianNB:} Asume que las características siguen una distribución normal (gaussiana) continua, siendo útil cuando las variables son continuas y tienen una forma aproximadamente normal.
\end{itemize}

\section{Variables Continuas: Gaussian Naive Bayes}
Hasta ahora vimos cómo trabajar con variables categóricas, donde simplemente contamos frecuencias. Pero, ¿qué pasa cuando tenemos variables continuas como edad, peso, altura, temperatura, etc.?

\subsection{Diferencia entre Variables Categóricas y Continuas}
\begin{itemize}
    \item \textbf{Variables categóricas:} Tienen valores discretos (Bajo, Medio, Alto). Podemos contar cuántas veces aparece cada valor.
    \item \textbf{Variables continuas:} Pueden tomar cualquier valor real (25.3, 70.5, 1.78). No podemos contar frecuencias porque cada valor puede ser único.
\end{itemize}

\subsection{Solución: Asumir Distribución Normal}
Para variables continuas, Gaussian Naive Bayes asume que los valores siguen una distribución normal (gaussiana). En lugar de contar frecuencias, calculamos dos parámetros para cada característica en cada clase:

\begin{itemize}
    \item \textbf{Media ($\mu$):} Promedio de los valores
    \item \textbf{Desviación estándar ($\sigma$):} Dispersión de los valores
\end{itemize}

\subsection{Fórmula de la Distribución Normal}
La probabilidad de un valor $x$ dado una clase se calcula con la función de densidad de probabilidad gaussiana:

\[
P(x|y) = \frac{1}{\sqrt{2\pi\sigma^2}} \exp\left(-\frac{(x-\mu)^2}{2\sigma^2}\right)
\]

donde:
\begin{itemize}
    \item $x$ es el valor de la característica que queremos evaluar
    \item $\mu$ es la media de esa característica para la clase $y$
    \item $\sigma$ es la desviación estándar de esa característica para la clase $y$
    \item $\pi \approx 3.14159$
    \item $e \approx 2.71828$ (base del logaritmo natural)
\end{itemize}

\textbf{Procedimiento para clasificar con Gaussian NB:}
\begin{enumerate}
    \item Para cada característica y cada clase, calcular $\mu$ y $\sigma$
    \item Para un nuevo dato, aplicar la fórmula gaussiana usando los parámetros de cada clase
    \item Multiplicar las probabilidades de todas las características (asumiendo independencia)
    \item Elegir la clase con mayor probabilidad posterior
\end{enumerate}

\section{Ejercicio 1: Clasificación de Estudiantes}
Para entender el funcionamiento del clasificador bayesiano ingenuo, usemos un caso de ejemplo: queremos construir un clasificador que prediga si un estudiante va a aprobar un examen, dado:

\begin{itemize}
    \item \textbf{Tiempo de estudio:} Bajo, Moderado o Alto
    \item \textbf{Método de estudio:} Lectura, Práctica, o Lectura y Práctica (L y P)
    \item \textbf{Puntuación en exámenes anteriores:} Bajo, Promedio o Alto
\end{itemize}

\begin{table}[H]
    \centering
    \begin{tabular}{cccc}
        \toprule
        Tiempo de estudio & Método de estudio & Puntuación & Resultado \\
        \midrule
        Bajo & Lectura & Bajo & Desaprobó \\
        Bajo & Práctica & Alta & Aprobó \\
        Moderado & Lectura y Práctica & Promedio & Aprobó \\
        Alto & Lectura y Práctica & Alta & Aprobó \\
        Alto & Lectura & Alta & Desaprobó \\
        Bajo & Lectura y Práctica & Baja & Desaprobó \\
        Alto & Práctica & Alta & Aprobó \\
        Moderado & Lectura & Alta & Aprobó \\
        Moderado & Lectura y Práctica & Promedio & Aprobó \\
        Moderado & Práctica & Bajo & Desaprobó \\
        \bottomrule
    \end{tabular}
\end{table}

\subsection{Tablas de Frecuencia}
A partir de los datos, construimos las tablas de frecuencia para cada atributo:

\vspace{1em}

\textbf{Tiempo de Estudio:}
\begin{table}[H]
    \centering
    \begin{tabular}{cccc}
        \toprule
        Tiempo de estudio & Aprobó & Desaprobó & Total \\
        \midrule
        Bajo & 1 & 2 & 3 \\
        Moderado & 3 & 1 & 4 \\
        Alto & 2 & 1 & 3 \\
        \midrule
        Total & 6 & 4 & 10 \\
        \bottomrule
    \end{tabular}
\end{table}

\vspace{1em}

\textbf{Método de Estudio:}
\begin{table}[H]
    \centering
    \begin{tabular}{cccc}
        \toprule
        Método de estudio & Aprobó & Desaprobó & Total \\
        \midrule
        Lectura & 1 & 2 & 3 \\
        Práctica & 2 & 1 & 3 \\
        L y P & 3 & 1 & 4 \\
        \midrule
        Total & 6 & 4 & 10 \\
        \bottomrule
    \end{tabular}
\end{table}

\vspace{1em}

\textbf{Puntuación en exámenes anteriores:}
\begin{table}[H]
    \centering
    \begin{tabular}{cccc}
        \toprule
        Puntuación & Aprobó & Desaprobó & Total \\
        \midrule
        Bajo & 0 & 3 & 3 \\
        Promedio & 2 & 0 & 2 \\
        Alto & 4 & 1 & 5 \\
        \midrule
        Total & 6 & 4 & 10 \\
        \bottomrule
    \end{tabular}
\end{table}

\subsection{Aplicación del Teorema de Bayes}
El clasificador bayesiano ingenuo asume que los atributos son independientes entre sí, por lo que podemos multiplicar las probabilidades. Para un nuevo estudiante, calculamos:

\begin{align*}
P(\text{Aprobó}) &= \frac{6}{10} = 0.6 \\
P(\text{Desaprobó}) &= \frac{4}{10} = 0.4
\end{align*}

Para clasificar un nuevo caso, multiplicamos las probabilidades condicionales por la probabilidad a priori. Por ejemplo, para un estudiante con:
\begin{itemize}
    \item Tiempo de estudio: Alto
    \item Método: Lectura y Práctica
    \item Puntuación previa: Alta
\end{itemize}

\begin{align*}
P(\text{Aprobó}|X) &= P(\text{Aprobó}) \times P(\text{Alto}|\text{Aprobó}) \\
&\quad \times P(\text{L y P}|\text{Aprobó}) \times P(\text{Alta}|\text{Aprobó}) \\
&= 0.6 \times \frac{2}{6} \times \frac{3}{6} \times \frac{4}{6} \\
&= 0.13
\end{align*}

\begin{align*}
P(\text{Desaprobó}|X) &= P(\text{Desaprobó}) \times P(\text{Alto}|\text{Desaprobó}) \\
&\quad \times P(\text{L y P}|\text{Desaprobó}) \times P(\text{Alta}|\text{Desaprobó}) \\
&= 0.4 \times \frac{1}{4} \times \frac{1}{4} \times \frac{1}{4} \\
&= 0.006
\end{align*}

Como $0.13 > 0.006$, el clasificador predice que el estudiante aprobará el examen.

\section{Ejercicio 2: Clasificación de Riesgo Crediticio}
Ahora vamos a trabajar con un ejemplo un poco más complejo que incluye tres clases y el problema de probabilidad cero que requiere suavizado de Laplace.

\subsection{Planteamiento del Problema}
Un banco quiere clasificar a sus clientes en tres categorías de riesgo crediticio: \textbf{Bajo}, \textbf{Medio} y \textbf{Alto}. Para esto, utiliza las siguientes características:

\begin{itemize}
    \item \textbf{Nivel de Ingreso:} Bajo, Medio, Alto
    \item \textbf{Historial Crediticio:} Malo, Regular, Bueno
    \item \textbf{Tipo de Empleo:} Temporal, Permanente
    \item \textbf{Nivel de Deudas:} Bajo, Alto
\end{itemize}

\subsection{Datos de Entrenamiento}
\begin{table}[H]
    \centering
    \small
    \begin{tabular}{ccccc}
        \toprule
        Ingreso & Historial & Empleo & Deudas & Riesgo \\
        \midrule
        Alto & Bueno & Permanente & Bajo & Bajo \\
        Alto & Bueno & Permanente & Bajo & Bajo \\
        Medio & Bueno & Permanente & Bajo & Bajo \\
        Alto & Regular & Permanente & Bajo & Bajo \\
        Medio & Bueno & Temporal & Bajo & Medio \\
        Medio & Regular & Permanente & Alto & Medio \\
        Bajo & Regular & Temporal & Alto & Medio \\
        Medio & Regular & Temporal & Bajo & Medio \\
        Bajo & Malo & Temporal & Alto & Alto \\
        Bajo & Malo & Temporal & Alto & Alto \\
        Bajo & Regular & Temporal & Alto & Alto \\
        Medio & Malo & Temporal & Alto & Alto \\
        Bajo & Malo & Permanente & Alto & Alto \\
        Alto & Malo & Temporal & Alto & Medio \\
        Medio & Bueno & Permanente & Alto & Medio \\
        \bottomrule
    \end{tabular}
\end{table}

\subsection{Tablas de Frecuencia}

\textbf{Nivel de Ingreso:}
\begin{table}[H]
    \centering
    \begin{tabular}{ccccc}
        \toprule
        Ingreso & Riesgo Bajo & Riesgo Medio & Riesgo Alto & Total \\
        \midrule
        Alto & 3 & 1 & 0 & 4 \\
        Medio & 1 & 4 & 1 & 6 \\
        Bajo & 0 & 1 & 4 & 5 \\
        \midrule
        Total & 4 & 6 & 5 & 15 \\
        \bottomrule
    \end{tabular}
\end{table}

\textbf{Historial Crediticio:}
\begin{table}[H]
    \centering
    \begin{tabular}{ccccc}
        \toprule
        Historial & Riesgo Bajo & Riesgo Medio & Riesgo Alto & Total \\
        \midrule
        Bueno & 3 & 2 & 0 & 5 \\
        Regular & 1 & 2 & 1 & 4 \\
        Malo & 0 & 2 & 4 & 6 \\
        \midrule
        Total & 4 & 6 & 5 & 15 \\
        \bottomrule
    \end{tabular}
\end{table}

\textbf{Tipo de Empleo:}
\begin{table}[H]
    \centering
    \begin{tabular}{ccccc}
        \toprule
        Empleo & Riesgo Bajo & Riesgo Medio & Riesgo Alto & Total \\
        \midrule
        Permanente & 4 & 2 & 1 & 7 \\
        Temporal & 0 & 4 & 4 & 8 \\
        \midrule
        Total & 4 & 6 & 5 & 15 \\
        \bottomrule
    \end{tabular}
\end{table}

\textbf{Nivel de Deudas:}
\begin{table}[H]
    \centering
    \begin{tabular}{ccccc}
        \toprule
        Deudas & Riesgo Bajo & Riesgo Medio & Riesgo Alto & Total \\
        \midrule
        Bajo & 4 & 2 & 0 & 6 \\
        Alto & 0 & 4 & 5 & 9 \\
        \midrule
        Total & 4 & 6 & 5 & 15 \\
        \bottomrule
    \end{tabular}
\end{table}

\subsection{Clasificación de un Nuevo Cliente}
Queremos clasificar un nuevo cliente con las siguientes características:
\begin{itemize}
    \item Ingreso: Medio
    \item Historial: Bueno
    \item Empleo: Temporal
    \item Deudas: Bajo
\end{itemize}

\textbf{Probabilidades a priori:}
\begin{align*}
P(\text{Riesgo Bajo}) &= \frac{4}{15} \approx 0.267 \\
P(\text{Riesgo Medio}) &= \frac{6}{15} = 0.4 \\
P(\text{Riesgo Alto}) &= \frac{5}{15} \approx 0.333
\end{align*}

\textbf{Cálculo para Riesgo Bajo (sin suavizado):}
\begin{align*}
P(\text{Bajo}|X) &= P(\text{Bajo}) \times P(\text{Medio}|\text{Bajo}) \times P(\text{Bueno}|\text{Bajo}) \\
&\quad \times P(\text{Temporal}|\text{Bajo}) \times P(\text{Bajo deudas}|\text{Bajo}) \\
&= \frac{4}{15} \times \frac{1}{4} \times \frac{3}{4} \times \frac{0}{4} \times \frac{4}{4} \\
&= \frac{4}{15} \times \frac{1}{4} \times \frac{3}{4} \times 0 \times 1 = 0
\end{align*}

\textbf{¡Problema!} La probabilidad es cero porque ningún cliente de Riesgo Bajo tiene empleo Temporal en nuestros datos de entrenamiento. Esto hace que toda la multiplicación sea cero, lo cual no es razonable.

\subsection{Solución: Suavizado de Laplace}
Para evitar probabilidades cero, aplicamos suavizado de Laplace con $\alpha = 1$:

\[
P(x_i|y) = \frac{count(x_i, y) + \alpha}{count(y) + \alpha \times n}
\]

donde $n$ es el número de valores posibles para esa característica.

\subsubsection{Recalculando con Laplace ($\alpha = 1$):}

Para la característica Empleo (que tiene 2 valores posibles: Permanente y Temporal):
\begin{align*}
P(\text{Temporal}|\text{Bajo}) &= \frac{0 + 1}{4 + 1 \times 2} = \frac{1}{6} \approx 0.167
\end{align*}

Ahora recalculamos todas las probabilidades con suavizado:

\vspace{0.5cm}

\textbf{Riesgo Bajo:}
\begin{align*}
P(\text{Medio}|\text{Bajo}) &= \frac{1 + 1}{4 + 1 \times 3} = \frac{2}{7} \approx 0.286 \\
P(\text{Bueno}|\text{Bajo}) &= \frac{3 + 1}{4 + 1 \times 3} = \frac{4}{7} \approx 0.571 \\
P(\text{Temporal}|\text{Bajo}) &= \frac{0 + 1}{4 + 1 \times 2} = \frac{1}{6} \approx 0.167 \\
P(\text{Bajo deudas}|\text{Bajo}) &= \frac{4 + 1}{4 + 1 \times 2} = \frac{5}{6} \approx 0.833
\end{align*}

\begin{align*}
P(\text{Bajo}|X) &= 0.267 \times 0.286 \times 0.571 \times 0.167 \times 0.833 \\
&\approx 0.00303
\end{align*}

\textbf{Riesgo Medio:}
\begin{align*}
P(\text{Medio}|\text{Medio}) &= \frac{4 + 1}{6 + 1 \times 3} = \frac{5}{9} \approx 0.556 \\
P(\text{Bueno}|\text{Medio}) &= \frac{2 + 1}{6 + 1 \times 3} = \frac{3}{9} \approx 0.333 \\
P(\text{Temporal}|\text{Medio}) &= \frac{4 + 1}{6 + 1 \times 2} = \frac{5}{8} = 0.625 \\
P(\text{Bajo deudas}|\text{Medio}) &= \frac{2 + 1}{6 + 1 \times 2} = \frac{3}{8} = 0.375
\end{align*}

\begin{align*}
P(\text{Medio}|X) &= 0.4 \times 0.556 \times 0.333 \times 0.625 \times 0.375 \\
&\approx 0.01736
\end{align*}

\textbf{Riesgo Alto:}
\begin{align*}
P(\text{Medio}|\text{Alto}) &= \frac{1 + 1}{5 + 1 \times 3} = \frac{2}{8} = 0.25 \\
P(\text{Bueno}|\text{Alto}) &= \frac{0 + 1}{5 + 1 \times 3} = \frac{1}{8} = 0.125 \\
P(\text{Temporal}|\text{Alto}) &= \frac{4 + 1}{5 + 1 \times 2} = \frac{5}{7} \approx 0.714 \\
P(\text{Bajo deudas}|\text{Alto}) &= \frac{0 + 1}{5 + 1 \times 2} = \frac{1}{7} \approx 0.143
\end{align*}

\begin{align*}
P(\text{Alto}|X) &= 0.333 \times 0.25 \times 0.125 \times 0.714 \times 0.143 \\
&\approx 0.00106
\end{align*}

\subsection{Decisión Final}
Comparando las probabilidades posteriores:
\begin{itemize}
    \item $P(\text{Riesgo Bajo}|X) \approx 0.00303$
    \item $P(\text{Riesgo Medio}|X) \approx 0.01736$ \textbf{← Mayor probabilidad}
    \item $P(\text{Riesgo Alto}|X) \approx 0.00106$
\end{itemize}

\textbf{Conclusión:} El cliente se clasifica como \textbf{Riesgo Medio}.

\vspace{0.5em}

\textbf{Nota importante:} Sin el suavizado de Laplace, habríamos obtenido probabilidad cero para Riesgo Bajo, lo cual habría sido un resultado erróneo. El suavizado nos permite manejar combinaciones de características que no aparecieron en los datos de entrenamiento.

\section{Ejercicio 3: Gaussian Naive Bayes - Clasificación de Actividad Física}
Ahora vamos a trabajar con un ejemplo que utiliza variables continuas, donde aplicaremos Gaussian Naive Bayes.

\subsection{Planteamiento del Problema}
Un centro de salud quiere predecir si una persona realiza actividad física regular basándose en características físicas y de estilo de vida. La variable objetivo es binaria: \textbf{Hace Ejercicio} (Sí/No).

Las características medidas son:
\begin{itemize}
    \item \textbf{Edad} (años)
    \item \textbf{Peso} (kg)
    \item \textbf{Horas de sueño} (horas por día)
\end{itemize}

\subsection{Datos de Entrenamiento}
\begin{table}[H]
    \centering
    \begin{tabular}{cccc}
        \toprule
        Edad & Peso & Horas de sueño & Hace Ejercicio \\
        \midrule
        22 & 65 & 7.5 & Sí \\
        25 & 70 & 8.0 & Sí \\
        23 & 68 & 7.0 & Sí \\
        28 & 72 & 7.5 & Sí \\
        24 & 66 & 8.5 & Sí \\
        26 & 69 & 7.8 & Sí \\
        45 & 85 & 6.0 & No \\
        50 & 90 & 5.5 & No \\
        48 & 88 & 6.2 & No \\
        52 & 92 & 5.8 & No \\
        49 & 87 & 6.5 & No \\
        47 & 86 & 5.9 & No \\
        \bottomrule
    \end{tabular}
\end{table}

\subsection{Cálculo de Parámetros (Media y Desviación Estándar)}
Para cada característica y cada clase, necesitamos calcular la media ($\mu$) y la desviación estándar ($\sigma$).

\subsubsection{Clase: Hace Ejercicio = Sí (6 muestras)}

\textit{Edad:}
\begin{align*}
\mu_{\text{edad, Sí}} &= \frac{22 + 25 + 23 + 28 + 24 + 26}{6} = \frac{148}{6} \approx 24.67 \text{ años} \\
\sigma_{\text{edad, Sí}} &= \sqrt{\frac{\sum (x_i - \mu)^2}{n}} \\
&= \sqrt{\frac{(22-24.67)^2 + (25-24.67)^2 + ... + (26-24.67)^2}{6}} \\
&= \sqrt{\frac{7.11 + 0.11 + 2.78 + 11.11 + 0.44 + 1.78}{6}} \\
&= \sqrt{\frac{23.33}{6}} \approx 1.97 \text{ años}
\end{align*}

\textit{Peso:}
\begin{align*}
\mu_{\text{peso, Sí}} &= \frac{65 + 70 + 68 + 72 + 66 + 69}{6} = \frac{410}{6} \approx 68.33 \text{ kg} \\
\sigma_{\text{peso, Sí}} &= \sqrt{\frac{11.11 + 2.78 + 0.11 + 13.44 + 5.44 + 0.44}{6}} \\
&= \sqrt{\frac{33.33}{6}} \approx 2.36 \text{ kg}
\end{align*}

\textit{Horas de sueño:}
\begin{align*}
\mu_{\text{sueño, Sí}} &= \frac{7.5 + 8.0 + 7.0 + 7.5 + 8.5 + 7.8}{6} = \frac{46.3}{6} \approx 7.72 \text{ horas} \\
\sigma_{\text{sueño, Sí}} &= \sqrt{\frac{0.048 + 0.078 + 0.518 + 0.048 + 0.608 + 0.006}{6}} \\
&= \sqrt{\frac{1.306}{6}} \approx 0.47 \text{ horas}
\end{align*}

\subsubsection{Clase: Hace Ejercicio = No (6 muestras)}

\textit{Edad:}
\begin{align*}
\mu_{\text{edad, No}} &= \frac{45 + 50 + 48 + 52 + 49 + 47}{6} = \frac{291}{6} = 48.5 \text{ años} \\
\sigma_{\text{edad, No}} &= \sqrt{\frac{12.25 + 2.25 + 0.25 + 12.25 + 0.25 + 2.25}{6}} \\
&= \sqrt{\frac{29.5}{6}} \approx 2.22 \text{ años}
\end{align*}

\textit{Peso:}
\begin{align*}
\mu_{\text{peso, No}} &= \frac{85 + 90 + 88 + 92 + 87 + 86}{6} = \frac{528}{6} = 88 \text{ kg} \\
\sigma_{\text{peso, No}} &= \sqrt{\frac{9 + 4 + 0 + 16 + 1 + 4}{6}} \\
&= \sqrt{\frac{34}{6}} \approx 2.38 \text{ kg}
\end{align*}

\textit{Horas de sueño:}
\begin{align*}
\mu_{\text{sueño, No}} &= \frac{6.0 + 5.5 + 6.2 + 5.8 + 6.5 + 5.9}{6} = \frac{35.9}{6} \approx 5.98 \text{ horas} \\
\sigma_{\text{sueño, No}} &= \sqrt{\frac{0.0004 + 0.2304 + 0.0484 + 0.0324 + 0.2704 + 0.0064}{6}} \\
&= \sqrt{\frac{0.5884}{6}} \approx 0.31 \text{ horas}
\end{align*}

\subsection{Resumen de Parámetros}
\begin{table}[H]
    \centering
    \begin{tabular}{llcc}
        \toprule
        \textbf{Característica} & \textbf{Clase} & $\boldsymbol{\mu}$ & $\boldsymbol{\sigma}$ \\
        \midrule
        Edad & Sí & 24.67 & 1.97 \\
        Edad & No & 48.5 & 2.22 \\
        \midrule
        Peso & Sí & 68.33 & 2.36 \\
        Peso & No & 88.0 & 2.38 \\
        \midrule
        Horas de sueño & Sí & 7.72 & 0.47 \\
        Horas de sueño & No & 5.98 & 0.31 \\
        \bottomrule
    \end{tabular}
\end{table}

\subsection{Clasificación de una Nueva Persona}
Queremos clasificar una nueva persona con las siguientes características:
\begin{itemize}
    \item Edad: 30 años
    \item Peso: 75 kg
    \item Horas de sueño: 7 horas
\end{itemize}

\textbf{Probabilidades a priori:}
\begin{align*}
P(\text{Sí}) &= \frac{6}{12} = 0.5 \\
P(\text{No}) &= \frac{6}{12} = 0.5
\end{align*}

\textbf{Aplicando la fórmula gaussiana:}

Recordemos que:
\[
P(x|y) = \frac{1}{\sqrt{2\pi\sigma^2}} \exp\left(-\frac{(x-\mu)^2}{2\sigma^2}\right)
\]

\subsubsection{Para la clase ``Hace Ejercicio = Sí'':}

\textit{Probabilidad de Edad = 30:}
\begin{align*}
P(30|\text{Sí}) &= \frac{1}{\sqrt{2\pi(1.97)^2}} \exp\left(-\frac{(30-24.67)^2}{2(1.97)^2}\right) \\
&= \frac{1}{\sqrt{24.37}} \exp\left(-\frac{28.39}{7.76}\right) \\
&= \frac{1}{4.94} \exp(-3.66) \\
&= 0.202 \times 0.026 \approx 0.0053
\end{align*}

\textit{Probabilidad de Peso = 75:}
\begin{align*}
P(75|\text{Sí}) &= \frac{1}{\sqrt{2\pi(2.36)^2}} \exp\left(-\frac{(75-68.33)^2}{2(2.36)^2}\right) \\
&= \frac{1}{\sqrt{35.03}} \exp\left(-\frac{44.49}{11.14}\right) \\
&= \frac{1}{5.92} \exp(-3.99) \\
&= 0.169 \times 0.018 \approx 0.0031
\end{align*}

\textit{Probabilidad de Sueño = 7:}
\begin{align*}
P(7|\text{Sí}) &= \frac{1}{\sqrt{2\pi(0.47)^2}} \exp\left(-\frac{(7-7.72)^2}{2(0.47)^2}\right) \\
&= \frac{1}{\sqrt{1.39}} \exp\left(-\frac{0.52}{0.44}\right) \\
&= \frac{1}{1.18} \exp(-1.18) \\
&= 0.847 \times 0.308 \approx 0.261
\end{align*}

\subsubsection{Probabilidad posterior para ``Sí'':}
\begin{align*}
P(\text{Sí}|X) &= P(\text{Sí}) \times P(30|\text{Sí}) \times P(75|\text{Sí}) \times P(7|\text{Sí}) \\
&= 0.5 \times 0.0053 \times 0.0031 \times 0.261 \\
&\approx 0.0000021
\end{align*}

\subsubsection{Para la clase ``Hace Ejercicio = No'':}

\textit{Probabilidad de Edad = 30:}
\begin{align*}
P(30|\text{No}) &= \frac{1}{\sqrt{2\pi(2.22)^2}} \exp\left(-\frac{(30-48.5)^2}{2(2.22)^2}\right) \\
&= \frac{1}{\sqrt{30.98}} \exp\left(-\frac{342.25}{9.86}\right) \\
&= \frac{1}{5.57} \exp(-34.71) \\
&= 0.180 \times 1.25 \times 10^{-15} \approx 2.25 \times 10^{-16}
\end{align*}

Dado que $P(30|\text{No})$ es extremadamente pequeño (la persona de 30 años está muy lejos de la media de edad de quienes no hacen ejercicio, que es 48.5 años), la probabilidad posterior para ``No'' será prácticamente cero.

\subsection{Decisión Final}
Aunque ambas probabilidades son pequeñas en valor absoluto, lo importante es compararlas:

\begin{itemize}
    \item $P(\text{Sí}|X) \approx 0.0000021$ \textbf{← Mucho mayor}
    \item $P(\text{No}|X) \approx 2.25 \times 10^{-16}$
\end{itemize}

\textbf{Conclusión:} La persona se clasifica como que \textbf{SÍ hace ejercicio}.

\vspace{0.5em}

\textbf{Interpretación:} Una persona de 30 años con 75 kg y 7 horas de sueño está mucho más cerca del perfil de las personas que hacen ejercicio en nuestros datos (jóvenes, peso moderado, buen sueño) que del perfil de quienes no hacen ejercicio (mayor edad, mayor peso, menos sueño).

\vspace{0.5em}

\textbf{Nota importante:} En Gaussian Naive Bayes, no necesitamos normalizar por $P(X)$ porque solo nos interesa comparar las probabilidades relativas. La clase con mayor probabilidad posterior no normalizada será siempre la misma que con normalización.

\section{Ventajas y Desventajas}
\underline{Ventajas}
\begin{itemize}
    \item Simple y rápido de entrenar
    \item Funciona bien con conjuntos de datos pequeños
    \item Maneja bien datos faltantes
    \item Fácil de interpretar
\end{itemize}

\underline{Desventajas}
\begin{itemize}
    \item La suposición de independencia raramente se cumple en la realidad
    \item Sensible a características irrelevantes
    \item No captura relaciones complejas entre características
    \item Puede tener problemas con probabilidades cero (requiere suavizado)
\end{itemize}

\section{Aplicaciones Comunes}

\begin{itemize}
    \item \textbf{Clasificación de spam:} Identificación de mails no deseados basada en el contenido y características del mensaje.
    
    \item \textbf{Análisis de sentimientos:} Clasificación de opiniones y comentarios en redes sociales, reseñas de productos o noticias como positivos, negativos o neutrales.
    
    \item \textbf{Diagnóstico médico:} Ayuda en la identificación de enfermedades basándose en síntomas y resultados de pruebas médicas.
    
    \item \textbf{Detección de fraude:} Identificación de transacciones sospechosas o actividades fraudulentas en sistemas financieros.
    
    \item \textbf{Clasificación de malware:} Identificación de software malicioso basado en sus características y comportamiento.
\end{itemize}

\section{Suavizado de Laplace}
Un problema común en Naive Bayes es cuando una probabilidad condicional es cero, lo que hace que toda la multiplicación sea cero. Para evitar esto, se utiliza el suavizado de Laplace (o suavizado aditivo), que agrega un pequeño valor $\alpha$ (típicamente 1) a todos los conteos. Los conteos representan la frecuencia con la que aparece cada valor de una característica para una clase específica.

\[
P(x_i|y) = \frac{count(x_i, y) + \alpha}{count(y) + \alpha \times n}
\]

Por ejemplo, si tenemos:
\begin{itemize}
    \item $\alpha = 1$ (suavizado de Laplace)
    \item $count(x_i, y) = 0$ (la palabra nunca apareció en la clase)
    \item $count(y) = 10$ (total de documentos en la clase)
    \item $n = 1000$ (cantidad de características)
\end{itemize}

Entonces:
\[
P(x_i|y) = \frac{0 + 1}{10 + 1 \times 1000} = \frac{1}{1010} \approx 0.001
\]

\vspace{0.5em}

Con esto, se evita que una probabilidad condicional sea cero, haciendo que toda la multiplicación sea cero.

\section{Implementación en Python}

\begin{lstlisting}[language=Python]
from sklearn.naive_bayes import GaussianNB
from sklearn.naive_bayes import MultinomialNB
from sklearn.naive_bayes import BernoulliNB

# Para datos continuos (distribucion normal)
model = GaussianNB()

# Para datos discretos (conteos)
model = MultinomialNB()

# Para datos binarios (presencia/ausencia)
model = BernoulliNB()

# Para datos categoricos
model = CategoricalNB()

# Para datos desbalanceados
model = ComplementNB()

# Entrenamiento
model.fit(X_train, y_train)

# Prediccion
y_pred = model.predict(X_test)
\end{lstlisting}

\end{document}
